\documentclass[10pt,a4paper]{amsart}
\usepackage[margin=19mm]{geometry}
\usepackage{amsmath,amssymb,mathtools}
\usepackage[scr=rsfs,cal=euler]{mathalfa}
\usepackage{xfrac}
\usepackage[unicode,hidelinks,bookmarks=false]{hyperref}
\newcommand{\ie}{i.e.\ }
\newcommand{\cf}{cf.\ }
\newcommand{\resp}{respectively\ }
\newcommand{\Subsection}{\subsection}
\providecommand{\nobreakdash}{\nobreak\hbox{-}}
\setlength{\emergencystretch}{2em}
\multlinegap0pt
\usepackage{bm}
\usepackage{amssymb}
\usepackage{dsfont}
\usepackage{tikz}
\newtheorem{proposition}{Proposition}
\newcommand{\Ge}[2]{\bm{\Gel(} #1,#2 \bm{)} }
\DeclareMathOperator{\Gel}{Ge}
\DeclareMathOperator{\Id}{Id}
\newcommand{\Res}{\mathcal{R}}
\newcommand{\Vityone}[1]{\Vert #1 \Vert_{L^{\infty}L^{1}}}
\newcommand{\Vvert}{\vert \hspace{-1pt} \vert \hspace{-1pt} \vert}
\title{Exponential Tail Estimates for Multitype Poisson Branching Processes and Application to Hawkes Processes}
\author{Th\'eo Leblanc}
\date{Source: arXiv:2507.08462v2 (CC BY 4.0)}
\begin{document}
\maketitle

\noindent\emph{Source and licence.} Exponential Tail Estimates for Multitype Poisson Branching Processes and Application to Hawkes Processes. Version arXiv:2507.08462v2 (CC BY 4.0), distributed by arXiv under the Creative Commons Attribution 4.0 International licence, http://creativecommons.org/licenses/by/4.0/, as recorded in the arXiv licence metadata for this exact version and shown on the abstract page https://arxiv.org/abs/2507.08462v2. Full licence notice: \texttt{licenses/arxiv-2507.08462-CC-BY-4.0.txt}.\par\smallskip

\noindent\textit{Editorial scope.} Counted unit: the article's proposition prop4.3 (source0.tex lines 1272--1283). The article's complete original proof of it is delivered with it (source0.tex lines 1346--1376, one \texttt{proof} environment of 31 lines, which the article places away from the statement). No mathematics is written, repaired or re-derived: the statement and the proof are the article's own lines, in the article's own notation, and no retained line is substituted or re-flowed.  Retained uncounted as the source-local context the proof needs: lines 1261--1264.  Nothing was omitted from any retained span, so the package carries no omission record.  No retained line contains a citation, so the wrapper carries no bibliography and no bibliography entry.  No retained line contains a figure environment, an image-inclusion command or drawing code, so no image is delivered and the package declares no asset dependency.  Editorial formatting only, in addition: an AMS \texttt{amsart} preamble that restates the article's own declarations and macros used by the retained lines, namely the environments \texttt{proposition} on separate counters, together with the standard packages those lines need.  The retained environments print in this document as Proposition 1 in order of appearance, whereas the article prints the same items with the numbering of its own full document.  The complete native source of the article as distributed in the arXiv e-print for this exact version is bundled unchanged as \texttt{sources/181414/source0.tex}.
\begin{proposition}\label{prop4.2}
	Let $\bm{v}$ be a nonnegative bivariate matrix-valued function. Suppose that $\bm{V} =\Vityone{\bm{v}}$ has spectral radius smaller than one, then the function $\bm{\psi} = \sum_{n\geq 1} \bm{v}^{\star n}$ is well defined and satisfies for all $d>0$ and all $p\in(0,1)$,
	\[\Res_{\bm{\psi}}^{\infty}(d) \preceq \sum_{k=0}^{\infty} \bm{V}^k \Res_{\bm{v}}^{\infty}\big((1-p)p^k d\big) (\Id-\bm{V})^{-1}. \]  
\end{proposition}

\begin{proposition}\label{prop4.3}
	Let $\bm{v}$ be a nonnegative bivariate matrix-valued functions. Suppose that $\bm{V} =\Vityone{\bm{v}} \in\Ge{\alpha}{K}$ with $\alpha<1$ and that
	\[\Res_{\bm{v}}^{\infty}(t) \preceq \frac{A}{(1+t)^{\gamma}} \bm{V}, \quad \forall t\geq 0.\]
	Then the function $\bm{\psi} = \sum_{n\geq 1} \bm{v}^{\star n}$ is well defined and satisfies, for all $t> 0$ and all $0<\delta <1$,
	\begin{equation*}
	\Res_{\bm{\psi}}^{\infty}(t) \preceq \frac{A \big(1+\frac{\gamma}{\delta \log(1/\alpha)}\big)^{\gamma}}{(1+t)^{\gamma(1-\delta)}}\bm{V} (\Id-\bm{V})^{-2} + \bm{V}^{\lfloor\frac{\gamma(1-\delta) \log(1+t)}{\log(1/\alpha)}\rfloor+2} (\Id-\bm{V})^{-2}.
	\end{equation*}
	In particular,
	\begin{equation*}
		\Vvert \Res_{\bm{\psi}}^{\infty}(t) \Vvert_{\infty} \leq \frac{\alpha K(K+1) \Big[ A\big(1+\frac{\gamma}{\delta \log(1/\alpha)}\big)^{\gamma} + 1\Big]}{(1-\alpha)^2(1+t)^{\gamma(1-\delta)}}, \quad \forall t>0.
	\end{equation*}
\end{proposition}

\begin{proof}[Proof of Proposition \ref{prop4.3}]
	Let $t > 0$. We use Proposition \ref{prop4.2} with $p = \exp(-\delta / \varepsilon)$ and cut the sum at $k_0 = \lfloor\varepsilon \log(1+ t)\rfloor \geq 0$ with \[\varepsilon = \frac{\gamma(1-\delta)}{\log(1/\alpha)}, \quad \text{and thus} \quad p = \exp(-\delta/\varepsilon) = \exp\Big(-\frac{\delta \log(1/\alpha)}{\gamma (1-\delta)} \Big). \]
	Thus we have
	\begin{align*}
		\Res^{\infty}_{\bm{\psi}}(t) & \preceq \sum_{k=0}^{k_0} \bm{V}^k \Res^{\infty}_{\bm{v}}((1-p)p^k t)(\Id-\bm{V})^{-1} + \sum_{k=k_0+1}^{\infty} \bm{V}^k \Res^{\infty}_{\bm{v}}((1-p)p^k t)(\Id-\bm{V})^{-1} \\
		& \preceq  (\Id-\bm{V})^{-1} \Res_{\bm{v}}((1-p)p^{k_0} t)(\Id-\bm{V})^{-1} + (\Id-\bm{V})^{-1} \bm{V}^{k_0+2} (\Id-\bm{V})^{-1}. \\
	\end{align*}
	By definition we have
	\[ (1-p)p^{k_0} t\geq (1-p)p^{\varepsilon\log(1+ t)} t = ( 1 - p) (1+t)^{-\delta} t.\]
	Thus we have
	\begin{align*}
		\Res^{\infty}_{\bm{v}}((1-p)p^{k_0} t) & \preceq \frac{A}{\Big[ 1+ ( 1 -p) (1+t)^{-\delta}t\Big]^{\gamma}} \bm{V} \\
		& \preceq \frac{A (1-p)^{-\gamma}}{( 1+ t)^{\gamma(1-\delta)}} \bm{V}
	\end{align*}
	where we used the elementary inequality \[\frac{1}{1+c(1+x)^{-a}x} \leq  \frac{c^{-1}}{(1+x)^{1-a}}, \quad x\geq 0, \ a,c\in (0,1). \]
	For $x>0$ we have $(1-e^{-x})^{-1}\leq 1+x^{-1}$ thus,
	\begin{align*}
		(1-p)^{-\gamma}  = \Big(1 - e^{-\frac{\delta \log(1/\alpha)}{\gamma(1-\delta)}}\Big)^{-\gamma}  \leq \Big(1 + \frac{\gamma(1-\delta)}{\delta \log(1/\alpha)}\Big)^{\gamma} & \leq \Big(1 + \frac{\gamma}{\delta \log(1/\alpha)}\Big)^{\gamma}.
	\end{align*}
	It proves the first result since $\bm{V}$ commutes with $(\Id-\bm{V})^{-1}$. Let us now check the second bound on the norm. Since we have for $j\geq 1$,
	\[ \Vvert \bm{V}^j(\Id-\bm{V})^{-1} \Vvert_{\infty} \leq \sum_{i\geq j} K\alpha^i = \frac{K \alpha^j}{1-\alpha},\]
	and that
	\[\Vvert (\Id-\bm{V})^{-1} \Vvert_{\infty} \leq 1+\sum_{i\geq 1} K\alpha^i \leq \frac{K +1}{1-\alpha},\]
	it is clear that 
	\begin{align*}
		\Vvert \Res^{\infty}_{\bm{\psi}}(t) \Vvert_{\infty}  & \leq \frac{K(K+1) \alpha}{(1-\alpha)^2} \times \frac{A \big(1 + \frac{\gamma}{\delta \log(1/\alpha)}\big)^{\gamma}}{( 1+ t)^{\gamma(1-\delta)}} + \frac{K(K+1) \alpha^{k_0+2}}{(1-\alpha)^2}.
	\end{align*} 
	Since we have
	\[ \alpha^{k_0+2} \leq \alpha \alpha^{\varepsilon \log(1+t)} = \alpha (1+t)^{-\gamma(1-\delta)},\]
	the bound is clear.
\end{proof}

\end{document}
